\section{The Model Families Screened and How Candidates Were Eliminated}
\label{sec:families}

Six model families are screened on the representation of Section~\ref{sec:representation}: the $k$-nearest-neighbour regressor and the local thin-plate interpolant of that section, Gaussian process regression~\citep{Rasmussen2005_GP} with a scaled squared-exponential kernel plus white noise, gradient-boosted trees (XGBoost~\citep{TODO_XGBoost}), the baseline multilayer perceptron, and a residual network of the same width and depth whose blocks add to identity skips. Each family runs in both units of prediction, one row at a time and curve-wise, where it learns the spline coefficients and the range of each curve from the curve parameters; the curve-wise interpolant is the interpolation between curves of the outline. The twelve candidates per dataset and target are eliminated by successive halving over the training rows: every candidate is scored through the harness of Section~\ref{sec:optimization} on $\knobFamiliesBudgetsFirst$ training rows, the best $\knobFamiliesKeepsFirst$ of each unit on the frozen folds go on to $\knobFamiliesBudgetsSecond$, and the best $\knobFamiliesKeepsSecond$ of each unit to $\knobFamiliesBudgetsThird$. Each unit is ranked apart: on $\knobFamiliesBudgetsFirst$ rows the neutron-star design keeps about a hundred of its curves, too few for a curve-wise base, which at full data was the best representation of Section~\ref{sec:representation}. A smaller budget keeps rows evenly spaced along every curve, at least $\knobFamiliesPerCurve$ per curve, and fewer whole curves when the rows do not suffice for all. Each candidate has a fixed budget, its default settings and one seed; a fit past $\knobFamiliesMinutes$ minutes or $\knobFamiliesMemoryGb$\,GB is the family's wall at that number of rows, and a pointwise Gaussian process whose kernel matrix would pass the memory budget is not run. \Cref{tab:families} reports the significant figures each candidate keeps on the folds at the most rows it reached; Appendix~\ref{app:families} describes each family, its cost and what to look for in its row.

\InputIfFileExists{98_families_tab_families}{}{}
\InputIfFileExists{98_families_fig_scaling}{}{}

The local thin-plate interpolant survives on every pair and leads three of the four; \cref{tab:families} gives each pair's figures. For every candidate of the last round the test curves agree with the folds within a quarter of a figure, so the ranking is not an artefact of the folds. The short list for the tuning of Section~\ref{sec:precision} therefore never holds more than three families per pair. Across the four pairs, the nearest curve is the bar to beat: the rows along a curve are dense and neighbouring curves differ little, so a plain lookup of the closest training row predicts the mass and the charge on the test curves more accurately than the baseline network, which leads only on the neutron-star folds; a surrogate earns its place by beating this lookup, or by matching it while staying smooth enough for the likelihood surface of a future MCMC inference.

\Pickref{fig:families-scaling}{}{ follows each candidate through the rounds. }The pointwise Gaussian process meets its wall on $\famNsMassGprPointwiseWallRows$ rows on all four pairs: a single fit runs past $\knobFamiliesMinutes$ minutes, as the cube of the rows foretells. Fitted curve-wise, on one row per curve, it takes a fraction of a second and ties for the lead on the black-hole charge. The curve-wise unit needs rows along every curve as well as curves: on the neutron stars it stays below zero figures up to $\knobFamiliesBudgetsSecond$ rows, about ten per curve, and leads only on $\knobFamiliesBudgetsThird$, about seventy per curve. At their fixed budget the networks keep at most $\famNsMassMlpPointwiseFigures$ figures, the perceptron on the neutron-star mass after a fit of $\famNsMassMlpPointwiseSeconds$ seconds, against $\famNsMassRbfPointwiseSeconds$ seconds for the interpolant; they reach the short list only on the neutron-star charge, where every family stays near one figure.

Between the two units of the interpolant the figures are close but the ripple is not: on the neutron-star mass the curve-wise residual ripples about thirty times more across $\beta$, and five times more across $\lambda$, than the pointwise one, since each curve is rebuilt on its own. Against the data ceilings of $\nsEdaCeilMFigures$ and $\nsEdaCeilDFigures$ figures, the neutron-star survivors stay about one figure short on both targets; the black-hole survivors stay about four figures short of $\bhEdaCeilMFigures$ and $\bhEdaCeilDFigures$.
\FloatBarrier

\todo{Set the survivors beside earlier neural surrogates~\citep{Liodis2024,Kaushikk2026,Liu2026}. $H_3$ is open: the screen never kept the perceptron and the residual network on the same pair, so their ripples do not compare here.}

\subsection{The Model Families of the Screen}
\label{app:families}

This appendix describes each candidate of the family screen of Section~\ref{sec:families}: what the model is, how its cost grows with the number of training rows $N$, and what to look for in its row of \cref{tab:families}. Every candidate reads the same design: the inputs $(\log_{10}\rhoc, \beta, \lambda)$ for neutron stars and $(\rh, \beta)$ for black holes, standardized on the rows it is fitted on, and the targets $M$ and $Y = \log_{10}(\Dch/M)$. Every candidate is fitted through the same harness, on the same curves and with one thread, so two rows of the table differ only in the model.

A pointwise model maps one row, a point on one curve, to its target; it sees every row as an independent sample and knows nothing of the curves, and it gains from every row. A curve-wise model first fits each training curve along its normalized coordinate $s \in [0, 1]$ as a least-squares cubic B-spline on eight knots, twelve numbers with the curve's start and end, learns them from the curve parameters alone, and rebuilds a row's curve to evaluate at its $s$; it is smooth along each curve by construction, its error is set by the number of curves, $\prepNsTrainCurves$ for neutron stars against $\prepBhTrainCurves$ for black holes, and when the curves are few, as at $\knobFamiliesBudgetsFirst$ rows, the predicted start and end miss and the error grows past the value itself.

\begin{table}[!htb]
\centering
\small
\caption{The candidates of the family screen: what each model is, how its cost grows with the number of training rows $N$, and what to look for in its row of \cref{tab:families}.}
\label{tab:families-models}
\begin{tabular}{@{}p{0.12\textwidth}p{0.34\textwidth}p{0.15\textwidth}p{0.28\textwidth}@{}}
\toprule
Family & What it is & Cost & What to look for \\
\midrule
Weighted $k$-nearest neighbours & Mean target of the $k = 8$ nearest training points in the standardized inputs, weighted by inverse distance~\citep{ScikitLearn2011}; nothing is fitted. & Tree build about $N \log N$, a query about $\log N$; no wall. & The floor any family must beat; it jumps where the neighbours change, a large ripple, and cannot extrapolate. \\
Local thin-plate RBF & The $\knobFamiliesNeighbours$ nearest points interpolated exactly by a thin-plate spline $\phi(r) = r^2 \log r$ plus a linear part~\citep{SciPy2020}; a constant part where the neighbours lie on one plane. & One system of about a hundred unknowns per query; grows with the queries, not with $N$. & Strong where the data are dense, the best model of Section~\ref{sec:representation}; no notion of noise, and it can overshoot at the edges. \\
Gaussian process & Scaled squared-exponential kernel plus white noise, $\sigma^2 \exp(-\lVert x - x' \rVert^2 / 2\ell^2) + \sigma_n^2 \delta_{xx'}$, set by the marginal likelihood~\citep{Rasmussen2005_GP,ScikitLearn2011}; the posterior mean with its variance. & About $N^3$ operations per optimizer step and $8 N^2$ bytes: pointwise it passes the $\knobFamiliesMemoryGb$\,GB budget past about $2.2 \times 10^4$ rows and is a wall; curve-wise, no wall. & Smooth everywhere, the strongest pointwise model on $10^3$ rows; does it keep its lead or meet its wall first. A noise at its lower bound means it interpolates; its variance is a free error bar. \\
XGBoost & Sum of regression trees, each fitted to the residual of the trees before it~\citep{TODO_XGBoost}, library defaults: one hundred trees of depth at most six; curve-wise, the twelve outputs together. & About $N$ operations per tree; no wall. & Piecewise constant, a staircase with a large ripple, a constant extrapolation; a reference for a model without smoothness. \\
Multilayer perceptron & $\knobNetDepth$ hidden layers of $\knobNetWidth$ ReLU units, AdamW~\citep{Loshchilov2017_AdamW} at a learning rate of $\knobNetLr$ on a cosine schedule, batches of $\knobNetBatchSize$ rows, the mean squared error of the standardized target; stops after $\knobNetPatience$ epochs without gain on $\knobNetValidFraction$ of the training curves, at most $\knobNetMaxEpochs$, and keeps its best epoch~\citep{PyTorch2019}. & About $N$ operations per epoch, linear in $N$; no wall, but minutes at $10^5$ rows. & Smooth and gains from every row, so likely to rise from round to round; trails the interpolants on $10^3$ rows. Its error is set by the optimization, so a fit that stops early tells as much as its figures. \\
Residual network & The perceptron of the same width and depth with residual blocks, $h \leftarrow h + \mathrm{ReLU}(W h + b)$~\citep{He2016_ResNet}, trained exactly as it. & As the perceptron. & Beside the perceptron it tests $H_3$: whether the skips give a smaller ripple at the same number of weights. \\
\bottomrule
\end{tabular}
\end{table}

Each entry of \cref{tab:families} is $-\log_{10}$ of the relative error at its 95th percentile, of $M$ or of $\Dch$ rebuilt from $Y$: two figures is an error of one percent on all but the worst five percent of the rows, a negative entry is an error larger than the value itself, and the data ceilings of Sections~\ref{sec:eda-ns} and~\ref{sec:eda-bh} bound what any model can reach. The table reports the frozen folds, whole training curves held out in turn, on which the halving ranks, and the test curves are scored too and stay held out from every choice, so a candidate far better on the folds than on the test curves has been chosen by the folds' noise. The subscript is the number of training rows of the last round the candidate reached, a candidate eliminated on $\knobFamiliesBudgetsFirst$ rows keeping its $\knobFamiliesBudgetsFirst$ entry, and a wall is a fit past $\knobFamiliesMinutes$ minutes or $\knobFamiliesMemoryGb$\,GB at that number of rows. Every run also records the ripple of its residual, the root mean square of its second divided difference across each curve parameter at fixed points along the curves: a smooth model keeps it near zero and a model that jumps between neighbouring curves does not, even when its figures are high; it is reported with the scaling of each family.

